\clearpage
\section{РЕАЛИЗАЦИЯ ПРОГРАММНОГО РЕШЕНИЯ}

\subsection{Выбор инструментов разработки}

Программная реализация экспериментального стенда выполнена на языке Python. Выбор данного языка обусловлен его фактическим статусом стандартного инструмента прикладного анализа данных и машинного обучения, а также наличием зрелой экосистемы библиотек для работы с временными рядами, построения торговых стратегий и обучения RL-агентов. Для хранения и преобразования табличных данных используются \textit{pandas} и \textit{NumPy}; для классических моделей машинного обучения — \textit{scikit-learn}; для получения котировок — \textit{yfinance}; для RL-среды и обучения агентов — \textit{Gymnasium} и \textit{stable-baselines3}; для визуализации — \textit{matplotlib} [15, 16, 17, 19, 20, 24].

Выбранный стек ориентирован на локальный запуск на ноутбуке класса MacBook M3 Pro с оперативной памятью 18 ГБ. Это наложило ограничение на архитектуру эксперимента: в основной программный контур были включены сравнительно легкие классические ML-модели и два RL-алгоритма, которые можно обучать без распределенной инфраструктуры. Такой выбор соответствует прикладной логике дипломной работы, где важна не демонстрация максимально тяжелой вычислительной схемы, а создание воспроизводимого и понятного исследовательского стенда.

\subsection{Структура программного проекта}

Программный проект организован модульно, что облегчает как сопровождение кода, так и связь между методической частью диплома и практической реализацией. Директории распределены по функциональному назначению: загрузка и подготовка данных, стратегии, модели, RL-компоненты, расчет метрик, визуализация и папка отчетов. Такая структура делает экспериментальный контур прозрачным и позволяет расширять его без перестройки базовой архитектуры.

\begin{table}[H]
\caption{Основные модули программного проекта}
\label{tab:project_modules}
\begin{center}
\small
\begin{tabularx}{\textwidth}{|p{4cm}|p{4cm}|X|}
\hline
\textbf{Модуль} & \textbf{Файл / каталог} & \textbf{Назначение} \\
\hline
Загрузка данных & \texttt{src/data/download.py} & Получение исторических котировок через \texttt{yfinance}, сохранение OHLCV-наборов в CSV \\
\hline
Генерация признаков & \texttt{src/data/features.py} & Вычисление доходностей, скользящих средних, RSI, MACD, Bollinger Bands и целевых переменных \\
\hline
Baseline-стратегии & \texttt{src/strategies/} & Реализация Buy \& Hold и Moving Average Crossover \\
\hline
ML-модели & \texttt{src/models/ml\_models.py} & Обучение классификаторов, оценка качества и конвертация прогнозов в торговые сигналы \\
\hline
RL-среда & \texttt{src/rl/trading\_env.py} & Торговая среда Gymnasium с действиями hold, buy, sell и reward по изменению портфеля \\
\hline
Обучение агентов & \texttt{src/rl/train\_dqn.py}, \texttt{src/rl/train\_ppo.py} & Подготовка сценариев обучения и тестирования RL-агентов \\
\hline
Финансовые метрики & \texttt{src/metrics/financial\_metrics.py} & Расчет доходности, волатильности, коэффициента Шарпа, просадки и других показателей \\
\hline
Визуализация & \texttt{src/visualization/plots.py} & Построение графиков цены с сигналами, equity curve, drawdown и сравнительных диаграмм \\
\hline
Оркестрация эксперимента & \texttt{main.py} & Запуск полного pipeline исследования и сохранение результатов в \texttt{reports/} \\
\hline
\end{tabularx}
\end{center}

\end{table}

Как видно из таблицы~\ref{tab:project_modules}, модульная структура соответствует логике, сформированной во второй главе. Каждый этап исследования имеет собственное программное представление, а главный сценарий запуска объединяет эти этапы в единый pipeline. Это особенно важно для дипломного проекта, поскольку позволяет описывать не абстрактную систему, а конкретное программное решение с четко заданными границами ответственности модулей. Входом полного контура служат дневные биржевые котировки в формате CSV или данные, получаемые напрямую через \textit{yfinance}, а выходом --- таблицы признаков, файлы с предсказаниями моделей, журналы действий RL-агентов, итоговые CSV/JSON с метриками и PNG-графики equity curve, сигналов и просадки.

Укрупненная структура проекта и примеры команд запуска приведены в приложении~А.

\subsection{Модуль загрузки и подготовки данных}

Модуль \texttt{src/data/download.py} реализует загрузку исторических котировок через сервис Yahoo Finance. На вход он принимает тикер, дату начала и дату окончания периода, а на выходе возвращает DataFrame с колонками \texttt{Date}, \texttt{Open}, \texttt{High}, \texttt{Low}, \texttt{Close}, \texttt{Adj Close}, \texttt{Volume}. В модуле предусмотрены валидация тикера, проверка периода, обработка пустого результата и сохранение данных в CSV. Такое решение делает модуль пригодным как для автономного запуска, так и для включения в общий pipeline эксперимента. Наличие отдельной функции сохранения облегчает повторное использование уже загруженных наборов данных и снижает зависимость последующих этапов от состояния сети [24].

На практике модуль решает сразу две задачи. Первая — получение стандартизированного OHLCV-датасета. Вторая — формирование воспроизводимого артефакта, на который затем могут опираться feature engineering, обучение моделей и построение графиков. За счет этого даже при повторных экспериментах исследователь работает с фиксированным файлом, а не с потенциально изменившимся источником данных.

\subsection{Модуль генерации признаков}

Модуль \texttt{src/data/features.py} преобразует OHLCV-набор в матрицу признаков, используемую как ML-моделями, так и RL-средой. Реализованные признаки включают простую и логарифмическую доходность, скользящие средние, локальную волатильность, RSI, MACD, сигнальную линию MACD и полосы Боллинджера. В этом же модуле формируются целевые переменные \texttt{target\_direction} и \texttt{target\_return}, после чего строки с \textit{NaN}, возникающие из-за оконных вычислений, удаляются. На выходе модуль возвращает DataFrame, содержащий исходные рыночные поля, набор признаков и целевые переменные [5, 7].

Отдельное значение имеет функция временного разбиения данных \texttt{split\_time\_series}, реализующая деление на обучающую, валидационную и тестовую части без перемешивания. Такое решение согласует программную реализацию с методологией исследования и исключает типичную для финансовых задач ошибку, связанную с утечкой информации из будущего. На уровне кода это обеспечивает прямую воспроизводимость экспериментальных сценариев, описанных в тексте диплома.

\subsection{Реализация baseline-стратегий}

В каталоге \texttt{src/strategies/} реализованы две benchmark-стратегии: \texttt{BuyAndHoldStrategy} и \texttt{MovingAverageCrossoverStrategy}. Первая формирует постоянную длинную позицию и отражает пассивное владение активом. Вторая вычисляет короткую и длинную скользящие средние и открывает позицию только в случае превышения короткой средней над длинной. Обе стратегии принимают на вход DataFrame с колонками \texttt{Date} и \texttt{Close}, а на выходе возвращают унифицированный DataFrame с колонками \texttt{Date}, \texttt{Close}, \texttt{position}, \texttt{daily\_return}, \texttt{strategy\_return} и \texttt{portfolio\_value}.

Единый формат результата является важным инженерным решением. Он позволяет использовать один и тот же модуль финансовых метрик и одни и те же функции визуализации для всех классов подходов. Благодаря этому baseline-стратегии становятся не отдельными учебными примерами, а полноценными участниками общего сравнительного контура.

\subsection{Реализация ML-моделей}

Модуль \texttt{src/models/ml\_models.py} содержит реестр базовых классификаторов: \texttt{LogisticRegression}, \texttt{RandomForestClassifier} и \texttt{GradientBoostingClassifier}. В модуле реализована общая процедура обучения, оценки и сохранения предсказаний на тестовом периоде. На вход она получает обучающую, валидационную и тестовую выборки, список колонок признаков и название целевой переменной. Для каждой модели вычисляются Accuracy, Precision, Recall, F1-score и ROC-AUC, после чего прогноз направления преобразуется в торговую стратегию с помощью функции \texttt{convert\_predictions\_to\_strategy} [13, 14, 15].

Такой подход демонстрирует принципиальный для работы переход от прогноза к торговому действию. ML-модель не оценивается исключительно как классификатор; ее предсказания сразу интерпретируются как торговые сигналы, а затем становятся входом для расчета \texttt{strategy\_return} и \texttt{portfolio\_value}. Тем самым практическая часть проекта полностью поддерживает исследовательскую идею о необходимости сопоставления прогностического качества с финансовым эффектом.

\subsection{Реализация торговой RL-среды}

Центральным элементом RL-части является модуль \texttt{src/rl/trading\_env.py}, реализующий кастомную среду \texttt{TradingEnv}, совместимую со стандартом Gymnasium и библиотекой stable-baselines3. Среда принимает DataFrame с признаками, список используемых feature columns, начальный капитал, транзакционные издержки и длину окна наблюдений. На выходе среда формирует вектор наблюдения, величину награды, признаки завершения эпизода и служебную информацию о текущем состоянии портфеля. Агенту предоставляется дискретное пространство действий: \texttt{hold}, \texttt{buy}, \texttt{sell} [19, 20].

Важной особенностью среды является наличие истории действий, позиций и стоимости портфеля. Это делает возможным не только обучение, но и последующий анализ поведения агента на тестовом интервале. Кроме того, среда спроектирована так, чтобы reward напрямую отражал изменение стоимости портфеля, что облегчает интерпретацию результата и связывает RL-постановку с финансовым смыслом задачи.

\subsection{Обучение RL-агентов DQN и PPO}

Для запуска обучения RL-агентов реализованы отдельные CLI-скрипты \texttt{src/rl/train\_dqn.py} и \texttt{src/rl/train\_ppo.py}. Оба скрипта загружают подготовленный CSV, выбирают признаки, делят данные по времени, создают \texttt{TradingEnv}, обучают соответствующую модель и сохраняют артефакты в папку \texttt{models\_rl/}. После обучения скрипты автоматически прогоняют агента на тестовом интервале, сохраняют действия, equity curve и рассчитывают финансовые метрики [19, 20]. В реализованной версии по умолчанию используется \texttt{window\_size = 30}, \texttt{initial\_cash = 10000}, \texttt{transaction\_cost = 0.001} и \texttt{timesteps = 10000}. Для DQN дополнительно заданы \texttt{learning\_rate = 1e-4}, \texttt{buffer\_size = 5000}, \texttt{learning\_starts = 500}, \texttt{batch\_size = 64}; для PPO --- \texttt{learning\_rate = 3e-4}, \texttt{n\_steps = 512}, \texttt{batch\_size = 64}, \texttt{gae\_lambda = 0.95}.

Параметры DQN подобраны умеренными, чтобы обучение не было слишком тяжелым для локального ноутбука: ограничен размер replay buffer, снижены значения \texttt{learning\_starts} и \texttt{gradient\_steps}, используется компактная частота обновления. PPO, в свою очередь, обучается на \texttt{MlpPolicy} с умеренным числом шагов и размером батча. Это согласуется с ограничениями дипломного исследования, где приоритет отдается воспроизводимости и понятной инженерной схеме, а не экстремальной оптимизации гиперпараметров.

\subsection{Модуль расчета финансовых метрик}

Модуль \texttt{src/metrics/financial\_metrics.py} реализует набор функций для оценки торговых стратегий: \texttt{total\_return}, \texttt{annualized\_return}, \texttt{annualized\_volatility}, \texttt{sharpe\_ratio}, \texttt{max\_drawdown}, \texttt{win\_rate}, \texttt{number\_of\_trades}, \texttt{calmar\_ratio}, а также агрегирующую функцию \texttt{calculate\_all\_metrics}. Входом служит унифицированный DataFrame с историей стоимости портфеля, доходностей стратегии и позиции [27, 28].

Реализация модуля с отдельными функциями для каждой метрики обеспечивает прозрачность вычислений и облегчает проверку корректности результатов. С инженерной точки зрения это также означает, что расчет показателей может использоваться повторно как для baseline-стратегий, так и для ML- и RL-подходов, не требуя отдельных ветвей логики.

\subsection{Модуль визуализации результатов}

Для интерпретации экспериментов создан модуль \texttt{src/visualization/plots.py}, содержащий функции построения графика цены с торговыми сигналами, сравнения equity curves, отображения drawdown и столбчатых диаграмм по выбранным метрикам. Все функции принимают путь сохранения и возвращают объект \texttt{matplotlib Figure}, что упрощает как автоматическое формирование отчетов, так и интерактивный анализ в ноутбуках.

С точки зрения дипломной работы модуль визуализации важен не только как вспомогательный инструмент, но и как средство объяснения результатов. Графики позволяют увидеть, где именно стратегия входила в рынок, как рос или снижался капитал, насколько глубокими были периоды просадки и в каких сценариях RL-агент отличался от baseline-решений. Воспроизводимость запуска обеспечивается единым CLI-интерфейсом, параметром \texttt{random\_state}, сохранением входных данных и записью промежуточных артефактов в каталог \texttt{reports/}.

\subsection{Выводы по главе}

В третьей главе было показано, как методологическая схема исследования реализуется в конкретном программном проекте на Python. Описана структура модулей, соответствующая этапам экспериментального контура, а также особенности реализации загрузки данных, генерации признаков, baseline-стратегий, ML-моделей, RL-среды, финансовых метрик и визуализации.

Таким образом, практическая часть работы представляет собой не набор разрозненных скриптов, а цельную исследовательскую систему. Это создает основу для четвертой главы, где рассматриваются результаты экспериментального моделирования и выполняется содержательное сравнение всех доступных подходов по единому набору показателей.
